\documentclass[12pt,a4paper,reqno]{amsart}

% language
\usepackage[greek,english]{babel}
\usepackage[utf8]{inputenc}
\usepackage{alphabeta}

% change default names to greek
\addto\captionsenglish{
    \renewcommand{\contentsname}{Περιεχόμενα}
    \renewcommand{\refname}{Βιβλιογραφία}
    \renewcommand{\datename}{Ημερομηνία:}
    \renewcommand{\urladdrname}{Ιστοσελίδα}
}

% math 
\usepackage{amsmath,amsthm,amssymb,amscd}

% font
\usepackage[cal=euler]{mathalfa}
\usepackage{libertinus-type1}
% \usepackage{txfonts} % for upright greek letters
\usepackage{bm} % for bold symbols
\usepackage{bbm} % for the simply-looking bb symbols

% miscellaneous 
\usepackage{changepage} %for indenting environments
\usepackage{csquotes} % example: \textcquote{}
\usepackage{blkarray}
\setcounter{MaxMatrixCols}{20} % default for pmatrix is 10!!
\usepackage{ytableau}
\usepackage{array} %needed to increase the vertical length in a tabular

% drawing
\usepackage{tikz,tikz-cd}
\usetikzlibrary{shapes.misc, patterns, matrix, calc, intersections,positioning,arrows.meta}
\usepackage{graphics,graphicx}
\usepackage{float} % provides enhanced control and customization options for floating objects such as figures and tables

% colors
\usepackage{xcolor}
\definecolor{darkcandyapplered}{rgb}{0.64, 0.0, 0.0}
\definecolor{midnightblue}{rgb}{0.1, 0.1, 0.44}
\definecolor{mylightblue}{HTML}{336699}
\definecolor{burntorange}{rgb}{0.8, 0.33, 0.0}
\definecolor{iceberg}{rgb}{0.44, 0.65, 0.82}
\definecolor{applegreen}{rgb}{0.55, 0.71, 0.0}
\definecolor{canaryyellow}{rgb}{1.0, 0.94, 0.0}
\definecolor{lightgray}{rgb}{0.83, 0.83, 0.83}

% hrefs
\usepackage{hyperref}
\usepackage[noabbrev,capitalize]{cleveref}
\hypersetup{
    pdftoolbar=true,        
    pdfmenubar=true,        
    pdffitwindow=false,     
    pdfstartview={FitH},  % fits the width of the page to the window
    pdftitle={},
    pdfauthor={},
    pdfsubject={},
    pdfkeywords={},
    pdfnewwindow=true,  % links in new window
    colorlinks=true,  % false: boxed links; true: colored links
    linkcolor=darkcandyapplered,   % color of internal links
    citecolor=midnightblue,  % color of links to bibliography
    urlcolor=cyan,  % color of external links
    linktocpage=true  % changes the links from the section body to the page number
    }

% geometry
\textwidth=16cm 
\textheight=21cm 
\hoffset=-55pt 
\footskip=25pt

% thm envs (you might need to change the path)
\theoremstyle{definition}
\newtheorem{theorem}{Θεώρημα}
\newtheorem{proposition}[theorem]{Πρόταση}
\newtheorem{lemma}[theorem]{Λήμμα}
\newtheorem{corollary}[theorem]{Πόρισμα}
\newtheorem{conjecture}[theorem]{Εικασία}
\newtheorem{problem}[theorem]{Πρόβλημα}
\newtheorem*{claim}{Ισχυρισμός}
\newtheorem{observation}[theorem]{Παρατήρηση}
\newtheorem{definition}[theorem]{Ορισμός}
\newtheorem{question}[theorem]{Ερώτημα}
\newtheorem*{questions}{Ερωτήματα}
\newtheorem*{que}{Ερώτημα}
\newtheorem{example}[theorem]{Παράδειγμα}
\newtheorem{exercise}{Άσκηση}
\newtheorem*{zorn}{Λήμμα του Zorn}
\newtheorem*{example_6.6}{Παράδειγμα~6.6}

\theoremstyle{remark}
\newtheorem*{remark}{Παρατήρηση}

% fixes the correct numbering of environments
\numberwithin{theorem}{section}
\numberwithin{exercise}{section}
\numberwithin{equation}{section}

% math ops 
% fields
\newcommand{\NN}{\mathbbmss{N}} 
\newcommand{\ZZ}{\mathbbmss{Z}} 
\newcommand{\QQ}{\mathbbmss{Q}} 
\newcommand{\RR}{\mathbbmss{R}} 
\newcommand{\CC}{\mathbbmss{C}} 
\newcommand{\KK}{\mathbbmss{K}} 
\newcommand{\FF}{\mathbbmss{F}} 
\newcommand{\HH}{\mathbbmss{H}} 

% symmetric group
\newcommand{\fS}{\mathfrak{S}}  

% calligraphic 
\newcommand{\aA}{\mathcal{A}} 
\newcommand{\bB}{\mathcal{B}}
\newcommand{\cC}{\mathcal{C}}
\newcommand{\dD}{\mathcal{D}}
\newcommand{\eE}{\mathcal{E}}
\newcommand{\fF}{\mathcal{F}}
\newcommand{\hH}{\mathcal{H}}
\newcommand{\iI}{\mathcal{I}}
\newcommand{\lL}{\mathcal{L}}
\newcommand{\oO}{\mathcal{O}}
\newcommand{\pP}{\mathcal{P}}
\newcommand{\qQ}{\mathcal{Q}}
\newcommand{\sS}{\mathcal{S}}
\newcommand{\mM}{\mathcal{M}}
\newcommand{\uU}{\mathcal{U}}

% bold
\newcommand{\bfa}{\mathbf{a}}
\newcommand{\bfe}{\mathbf{e}}
\newcommand{\bfF}{\pmb{F}}
\newcommand{\bfR}{\pmb{R}}
\newcommand{\bfv}{\mathbf{v}}
%\newcommand{\bfx}{\bm{x}}
%\newcommand{\bfx}{\mathbf{x}} 
\newcommand{\bfx}{\pmb{x}}
\newcommand{\bfX}{\pmb{X}}
\newcommand{\bfy}{\pmb{y}}
\newcommand{\bfz}{\pmb{z}}

% roman
\newcommand{\rmA}{\mathrm{A}}
\newcommand{\rmB}{\mathrm{B}}
\newcommand{\rmC}{\mathrm{C}}
\newcommand{\rmD}{\mathrm{D}} 
\newcommand{\rmF}{\mathrm{F}} 
\newcommand{\rmH}{\mathrm{H}} 
\newcommand{\rmh}{\mathrm{h}} 
\newcommand{\rmI}{\mathrm{I}} 
\newcommand{\rmK}{\mathrm{K}}
\newcommand{\rmM}{\mathrm{M}}
\newcommand{\rmP}{\mathrm{P}}  
\newcommand{\rmp}{\mathrm{p}}  
\newcommand{\rmQ}{\mathrm{Q}}  
\newcommand{\rmR}{\mathrm{R}}
\newcommand{\rmS}{\mathrm{S}}
\newcommand{\rmT}{\mathrm{T}}
\newcommand{\rmU}{\mathrm{U}}
\newcommand{\rmV}{\mathrm{V}}
\newcommand{\rmY}{\mathrm{Y}}
\newcommand{\rmZ}{\mathrm{Z}}
\newcommand{\rmz}{\mathrm{z}}

% arrows and symbols 
\renewcommand{\to}{\rightarrow}
\newcommand{\toto}{\longrightarrow}
\newcommand{\mapstoto}{\longmapsto}
\newcommand{\then}{\Rightarrow}
\newcommand{\IFF}{\Leftrightarrow}
\newcommand{\tl}{\tilde}
\newcommand{\wtl}{\widetilde}
\newcommand{\ol}{\overline}
\newcommand{\ul}{\underline}
\newcommand{\oldemptyset}{\emptyset}
\renewcommand{\emptyset}{\varnothing}
\DeclareMathSymbol{\Arg}{\mathbin}{AMSa}{"39} % for arguments 
\newcommand{\onto}{\ensuremath{\twoheadrightarrow}}
\newcommand{\tle}{\trianglelefteq}
\newcommand{\tge}{\trianglerighteq}

% custom ops
\newcommand{\rmKer}{\mathrm{Ker}}
\newcommand{\rmIm}{\mathrm{Im}}
\newcommand{\lcm}{\mathrm{lcm}}
\newcommand{\GL}{\mathrm{GL}}
\newcommand{\ev}{\mathrm{ev}}
\newcommand{\End}{\mathrm{End}}
\newcommand{\rmid}{\mathrm{id}}
\newcommand{\rmspan}{\mathrm{span}}
\newcommand{\Hom}{\mathrm{Hom}}
\newcommand{\Ann}{\mathrm{Ann}}
\newcommand{\Set}{\pmb{\mathrm{Set}}}
\newcommand{\Rmod}{\pmb{\mathrm{mod}}}
\newcommand{\rmrank}{\mathrm{rank}}
\newcommand{\mSpec}{\mathrm{mSpec}}
\newcommand{\Spec}{\mathrm{Spec}}

% absolute value symbol
\usepackage{mathtools} 
\DeclarePairedDelimiter\abs{\lvert}{\rvert}%
\DeclarePairedDelimiter\norm{\lVert}{\rVert}%
\makeatletter
\let\oldabs\abs
\def\abs{\@ifstar{\oldabs}{\oldabs*}}

% tensor symbol
\newcommand{\tensor}[1]{%
  \mathbin{\mathop{\otimes}\limits_{#1}}%
}

% permutation cycle notation
\ExplSyntaxOn
\NewDocumentCommand{\cycle}{ O{\;} m }
 {
  (
  \alec_cycle:nn { #1 } { #2 }
  )
 }

\seq_new:N \l_alec_cycle_seq
\cs_new_protected:Npn \alec_cycle:nn #1 #2
 {
  \seq_set_split:Nnn \l_alec_cycle_seq { , } { #2 }
  \seq_use:Nn \l_alec_cycle_seq { #1 }
 }
\ExplSyntaxOff

% setminus symbol
\newcommand{\mysetminusD}{\hbox{\tikz{\draw[line width=0.6pt,line cap=round] (3pt,0) -- (0,6pt);}}}
\newcommand{\mysetminusT}{\mysetminusD}
\newcommand{\mysetminusS}{\hbox{\tikz{\draw[line width=0.45pt,line cap=round] (2pt,0) -- (0,4pt);}}}
\newcommand{\mysetminusSS}{\hbox{\tikz{\draw[line width=0.4pt,line cap=round] (1.5pt,0) -- (0,3pt);}}}
\newcommand{\sm}{\mathbin{\mathchoice{\mysetminusD}{\mysetminusT}{\mysetminusS}{\mysetminusSS}}}

% miscellaneous commands
\newcommand{\defn}[1]{{\color{mylightblue}{#1}}}
\newcommand{\toDo}{{\bf\color{red} TODO}}
\newcommand{\toCite}{{\bf\color{green} CITE}}
\newcommand*{\vertbar}{\rule[-1ex]{0.5pt}{2.5ex}} % for matrices with column vectors
\newcommand*{\horzbar}{\rule[.5ex]{2.5ex}{0.5pt}} % for matrices with row vectors
\newcommand{\myblue}[1]{{\color{iceberg}{#1}}}
\newcommand{\myorange}[1]{{\color{burntorange}{#1}}}
\newcommand{\mygreen}[1]{{\color{applegreen}{#1}}}
\newcommand{\myred}[1]{{\color{darkcandyapplered}{#1}}}

% ferrer's diagram
\newcommand{\fdiagram}[1]{
    \begin{tikzpicture}[scale=.7]
        \fill foreach \Z [count=\Y] in {#1}
        {foreach \X in {1,...,\Z} 
        {(\X,-\Y) circle[radius=3pt]}};
    \end{tikzpicture}
}

%
\usepackage{pgfplots}
\pgfplotsset{compat=1.18}

%
\makeatletter
\def\mathcolor#1#{\@mathcolor{#1}}
\def\@mathcolor#1#2#3{%
  \protect\leavevmode
  \begingroup
    \color#1{#2}#3%
  \endgroup
}
\makeatother

% 
\newenvironment{nouppercase}{%
  \let\uppercase\relax%
  \renewcommand{\uppercasenonmath}[1]{}}{}

% titlepage
\title{226: Θεωρία Δακτυλίων και Modules}
\author[Β.~Δ. Μουστακας]{Βασίλης Διονύσης Μουστάκας \\ Πανεπιστήμιο Κρήτης}
\date{31 Μαρτίου 2026}
% \urladdr{\href{https://sites.google.com/view/vasmous}{https://sites.google.com/view/vasmous}}

\begin{document}

\begingroup
\def\uppercasenonmath#1{} % this disables uppercase title
\let\MakeUppercase\relax % this disables uppercase authors
\maketitle
\endgroup

\setcounter{section}{6}
\setcounter{theorem}{11}
\setcounter{equation}{7}
\begin{center}
    \textbf{6. Τανυστικό γινόμενο
} (Συνέχεια)
\end{center}

Υπενθυμίζουμε πως σε ότι ακολουθεί, $R$ είναι ένας μεταθετικός δακτύλιος.

Έχοντας ορίσει το τανυστικό γινόμενο ενός αυθαίρετου πεπερασμένου πλήθους προτύπων, μπορούμε να ορίσουμε μία άλγεβρα, ξεκινώντας από ένα πρότυπο η οποία θα το περιέχει και θα είναι \emph{καθολική} ως προς αυτή την ιδιότητα, όπως κάναμε διαισθητικά και στην περίπτωση του πολυωνυμικού δακτύλιου.

\begin{definition}
    \label{def:tensor_algebra}
    Έστω $M$ ένα $R$-πρότυπο. Για κάθε $k \ge 0$, ορίζουμε 
    \[
    \rmT^k(M) \coloneqq M^{\otimes{k}} \coloneqq \underbrace{M\otimes M \otimes \cdots \otimes M}_{\text{$k$ φορές}},
    \]
    όπου $\rmT^0(M) \coloneqq R$ και θέτουμε 
    \[
    \rmT(M) \coloneqq \bigoplus_{k \ge 0} \rmT^k(M).
    \]
    Το $\rmT(M)$ ονομάζεται τανυστική άλγεβρα (tensor algebra) του $M$.
\end{definition}

\begin{example}
    \label{ex:tensor_algebra}
    Έστω $M$ ένα ελεύθερο $R$-πρότυπο με βάση $\{x_1, x_2, \dots, x_n\}$. Από την Πρόταση 6.7 το σύνολο 
    % \ref{prop:tensorProduct_free} 
    \[
    \{
        x_{i_1} \otimes x_{i_2} \otimes \cdots \otimes x_{i_k} : 1 \le i_j \le n, \ 1 \le j \le k
    \}
    \]
    αποτελεί βάση του $\rmT^k(M)$. Συνεπώς, το $\rmT^k(M)$ είναι ελεύθερο $R$-πρότυπο τάξης $n^k$. 

    Μπορούμε να φανταστούμε τα στοιχεία της τανυστικής άλγεβρας του $M$ ως $R$-γραμμι\-κούς συνδυασμούς \textquote{λέξεων} στο \textquote{αλφάβητο} $\{x_1, x_2, \dots, x_n\}$ με τον πολλαπλασιασμό να είναι η \emph{παράθεση} (concatenation) λέξεων. Ισοδύναμα το $\rmT(M)$ είναι ένας δακτύλιος \textquote{πολυω\-νύμων} στις μη μεταθετικές μεταβλητές $x_1, x_2, \dots, x_n$.
    
    Για παράδειγμα, για $n=2$ και $M = R\{x,y\}$, έχουμε 
    \[
    \rmT(M) = R \oplus R\{x, y\} \oplus R\{x\otimes x, x\otimes y, y \otimes x, y \otimes y\} \oplus \cdots. 
    \]
    Ένα παράδειγμα στοιχείου στο $\rmT(M)$ είναι 
    \[
    x \otimes x \otimes y - y \otimes x \otimes x + x \otimes y \otimes x,
    \]
    τ' οποίο θα μπορούσαμε να γράψουμε ως 
    \[
    x^2y - yx^2 + xyx.
    \]
\end{example}

\begin{theorem}
    \label{thm:tensor_algebra}
    Έστω $M$ ένα $R$-πρότυπο. Η τανυστική άλγεβρα $\rmT(M)$ με πολλαπλασιασμό
    \[
    (a_1 \otimes \cdots a_i)(a_1' \otimes \cdots a_j') = a_1 \otimes \cdots \otimes a_i \otimes a_1' \otimes \cdots a_j'
    \]
    είναι $R$-άλγεβρα και ικανοποιεί την εξής καθολική ιδιότητα: Για κάθε $R$-άλγεβρα $A$ και κάθε $R$-γραμμική απεικόνιση $\phi : M \to A$, υπάρχει μοναδικός ομομορφισμός $R$-αλγεβρών $\tilde{\phi} : \rmT(M) \to A$ τέτοιος ώστε το διάγραμμα 
    \[
    \begin{tikzcd}
    M \arrow[r, "\iota"] \arrow[dr, "\phi"'] & \rmT(M) \arrow[d, "\tilde{\phi}"] \\
                                        & A
    \end{tikzcd}
    \]
    να είναι μεταθετικό, ή ισοδύναμα $\tilde{\phi}(m) = \phi(m)$, για κάθε $m \in M$.
\end{theorem}

% \begin{proof}[Απόδειξη]    
% \end{proof}

\begin{example}
    \label{ex:tensor_algebra_2}
    Έχουμε $\rmT(\ZZ_n) \cong \ZZ[x]/(nx)$. Πράγματι, από το Παράδειγμα \ref{ex:tensorProduct} (δ) έπεται ότι 
    \[
    \rmT^k(\ZZ_n) = \ZZ_n^{\otimes{k}} = \ZZ_n.
    \]
    Επομένως, 
    \[
    \rmT(\ZZ_n) = \ZZ \oplus \ZZ_n \oplus \ZZ_n \oplus \cdots.
    \]
    Από την άλλη μεριά, στο πηλίκο $\ZZ[x]/(nx)$ μηδενίζονται τα στοιχεία του $\ZZ[x]$ της μορφής 
    \[
    nx, \ nx^2, \dots, \ nx^d, \ \dots.
    \]
    Επομένως, τα στοιχεία του $\ZZ[x]/(nx)$ είναι της μορφής 
    \[
    a_0 + a_1x + a_2x^2 + \cdots + a_dx^d
    \]
    όπου $a_0 \in \ZZ$ και $a_1, a_2, \dots, a_d \in \ZZ_n$. Με άλλα λόγια, 
    \begin{align*}
    \frac{\ZZ[x]}{(nx)} &= \frac{\ZZ \oplus \ZZ{x} \oplus \ZZ{x^2} \oplus \cdots}{(nx)} \\ 
    &\cong \ZZ \oplus \ZZ_n{x} \oplus \ZZ_n{x^2} \oplus \cdots \\
    &\cong \ZZ \oplus \ZZ_n \oplus \ZZ_n \oplus \cdots \\ 
    &= \rmT(\ZZ_n).
    \end{align*}
\end{example}

Το Θεώρημα \ref{thm:tensor_algebra} μας πληροφορεί ότι η τανυστική άλγεβρα $\rmT(M)$ είναι ο πιο \textquote{οικονομικός} τρόπος να μετατρέψουμε το $M$ σε $R$-άλγεβρα. Για τον λόγο αυτό, αν το $M$ είναι ελεύθερο πρότυπο με βάση $\{x_1, x_2, \dots, x_n\}$, η $\rmT(M)$ ονομάζεται \emph{ελεύθερη άλγεβρα} που παράγεται από τα $x_1, x_2, \dots, x_n$. Η τανυστική άλγεβρα δεν είναι (κατ' ανάγκη) μεταθετική.

\begin{definition}
    \label{def:symmetric_algebra}
    Έστω $M$ ένα $R$-πρότυπο. Το 
    \[
    \rmS(M) \coloneqq \frac{\rmT(M)}{R\{m \otimes m' - m' \otimes m : m, m' \in M\}}
    \]
    ονομάζεται \defn{συμμετρική άλγεβρα} (symmetric algebra) του $M$.
\end{definition}

\begin{example}
    \label{ex:symmetric_algebra}
    Έστω $M$ ένα ελεύθερο $R$-πρότυπο με βάση $\{x, y\}$. Στο Παράδειγμα \ref{ex:tensor_algebra} είδαμε ότι 
    \[
    \rmT(M) = R \oplus R\{x, y\} \oplus R\{x\otimes x, x\otimes y, y \otimes x, y \otimes y\} \oplus \cdots. 
    \]
    Στην περίπτωση αυτή, το υποπρότυπο που παράγεται από όλα τα $m \otimes m' - m' \otimes m$ είναι το 
    \[
    (x \otimes y - y \otimes x)
    \]
    (γιατί;). Κατά συνέπεια, περνώντας από την $\rmT(M)$ (την οποία σκεφτόμαστε ως ελεύθερη άλγεβρα) στην συμμετρική άλγεβρα $\rmS(M)$ ουσιαστικά κάνουμε την ταύτιση 
    \[
    x \otimes y = y \otimes x.
    \]
    Με άλλα λόγια, κάνουμε τις \textquote{μεταβλητές} $x$ και $y$ να μετατίθενται. Κατά συνέπεια, 
    \begin{align*}
    \rmS(M) &= R \oplus R\{x, y\} \oplus R\{x\otimes x, x \otimes y, y \otimes y\} \oplus \cdots \\
    &\cong R \oplus R\{x,y\} \oplus R\{x^2, xy, y^2\} \oplus \cdots \\ 
    &= R[x,y]. 
    \end{align*}
    Γενικότερα, αν το $M$ είναι ελεύθερο $R$-πρότυπο τάξης $n$, τότε η συμμετρική άλγεβρα του $M$ ταυτίζεται με τον πολυωνυμικό δακτύλιο σε $n$ μεταβλητές με συντελεστές από το $R$. Αυτή η ταύτιση είναι ένα παράδειγμα κατασκευής που \emph{δεν} εξαρτάται από τη βάση του $M$ και είναι ένας ακόμη λόγος που φανταζόμαστε τα στοιχεία $x_1, x_2, \dots, x_n$ ως μεταβλητές.
\end{example}

Κλείνουμε αυτή την παράγραφο με την καθολική ιδιότητα της συμμετρικής άλγεβρας, η οποία γενικεύει την καθολική ιδιότητα του πολυωνυμικού δακτύλιου που είδαμε στην Ενότητα 3.

\begin{theorem}
    \label{thm:symmetric_algebra}
    Έστω $M$ ένα $R$-πρότυπο. Η $\rmS(M)$ είναι μεταθετική άλγεβρα και ικανοποιεί την εξής καθολική ιδιότητα: Για κάθε μεταθετική $R$-άλγεβρα $A$ και κάθε $R$-γραμμική απεικό\-νιση $\phi : M \to A$, υπάρχει μοναδικός ομομορφισμός $R$-αλγεβρών $\tilde{\phi} : \rmS(M) \to A$ τέτοιος ώστε το διάγραμμα 
    \[
    \begin{tikzcd}
    M \arrow[r, "\iota"] \arrow[dr, "\phi"'] & \rmS(M) \arrow[d, "\tilde{\phi}"] \\
                                        & A
    \end{tikzcd}
    \]
    να είναι μεταθετικό, ή ισοδύναμα $\tilde{\phi}(m) = \phi(m)$, για κάθε $m \in M$.
\end{theorem}


% \begin{thebibliography}{99}
%     \bibitem{Alu09}
%     P.~Aluffi,
%     \textit{Algebra, Chapter 0},
%     Grad. Stud. in Math. {\bf104},
%     Amer. Math. Soc., Providence, RI, 2009.
%     \bibitem{DF04}
%     D. S. Dummit και R. M. Foote, 
%     \textit{Abstract algebra},
%     third edition, Wiley, 2004. 
% \end{thebibliography}
\end{document}